\documentclass[a4paper]{article}

\usepackage[a4paper, top=8mm, bottom=8mm, left=8mm, right=8mm]{geometry}

\usepackage{polyglossia}
\setdefaultlanguage[babelshorthands=true]{russian}
\setotherlanguage{english}

\usepackage{fontspec}
\setmainfont{FreeSerif}
\newfontfamily{\russianfonttt}[Scale=0.7]{DejaVuSansMono}

\usepackage[tiny, compact]{titlesec}

\usepackage{titling}
\setlength{\droptitle}{-1cm}
\pretitle{\begin{center}\begin{bfseries}\Large}
\posttitle{\par\end{bfseries}\end{center}}
\preauthor{\begin{center}\normalsize}
\postauthor{\par\end{center}\vspace{-1.8cm}}

\usepackage{hyperref}
\usepackage{bookmark}
\usepackage{csquotes}
\usepackage{booktabs}
\setlength{\emergencystretch}{1.5em}

\title{Spla: SSSP на GPU от Imagination Technologies}

\author{Редькин Алексей Андреевич}

\date{}

\begin{document}

\maketitle

\begin{flushright}
    Группа: \emph{2025.Б72-мм}

    Кафедра: \emph{Кафедра системного программирования}

    Научный руководитель: \emph{Григорьев Семён Вячеславович,}\\
    \emph{доцент кафедры системного программирования,}\\
    \emph{кандидат физико-математических наук}

    Консультант: \emph{Кутуев Владимир Александрович}

    Старший инженер-исследователь YADRO

    Номер семестра практики: \emph{2}
\end{flushright}

\section{Введение}

Библиотека Spla\footnote{URL: \url{https://github.com/SparseLinearAlgebra/spla} (дата обращения:~30.09.2026).}~--- это библиотека обобщённой разреженной линейной алгебры с открытым исходным кодом для математических вычислений с ускорением на графических процессорах (GPU): матрицы, векторы, скаляры и параметризуемые операции над ними. Использование открытого стандарта OpenCL для написания программ, выполняющих параллельные вычисления на различных устройствах, позволяет Spla работать на GPU различных производителей. Графовые алгоритмы реализуются поверх этих примитивов. Объектом исследования выступает плата Milk-V Jupiter\footnote{URL: \url{https://milkv.io/jupiter} (дата обращения:~07.10.2026).} на открытой архитектуре системы команд RISC-V с графическим процессором PowerVR (линейка графических процессоров компании Imagination Technologies)\footnote{URL: \url{https://www.imaginationtech.com/product/img-bxe-2-32/} (дата обращения:~07.10.2026).}, на которой обработка отдельных разреженных графов завершалась отказом алгоритма поиска кратчайших путей из одной вершины (Single-Source Shortest Paths, SSSP).

\section{Постановка задачи}

На целевом устройстве алгоритм поиска кратчайших путей из одной вершины преждевременно завершал работу на отдельных графах: значения расстояний не обновлялись и оставались равными \texttt{FLT\_MAX} (максимальное конечное значение типа \texttt{float}~--- маркер недостижимости вершины), а сверка с эталонной реализацией на CPU (центральный процессор) показывала расхождение более чем на 99\%~вершин на примере \texttt{belgium\_osm}.

\textbf{Целью работы является} локализация причины отказа и разработка переносимого решения для её устранения.

Для достижения цели были решены следующие задачи:
\begin{itemize}
    \item выполнить обзор итеративной схемы SSSP в Spla;
    \item локализовать место отказа в цепочке итераций алгоритма SSSP и экспериментально подтвердить корневую причину на целевом устройстве;
    \item реализовать переносимое исправление;
    \item провести проверку исправления сверкой с эталоном на наборе тестовых графов.
\end{itemize}

Результаты работы необходимы пользователям вычислительной платформы RISC-V с GPU PowerVR, для которых библиотека Spla~--- инструмент выполнения графовых алгоритмов с аппаратным ускорением. Корректная обработка некоторых разреженных графов на данной платформе не была гарантирована: на ряде графов алгоритм завершался отказом. В выполнении работы заинтересованы разработчики библиотеки Spla.

\section{Описание предлагаемого решения}

Замена в операции \texttt{MIN\_FLOAT} (минимум для чисел с плавающей точкой) функции \texttt{min} языка OpenCL C на \texttt{fmin} приводила к успешному завершению алгоритма, поэтому подозреваемой стала операция \texttt{MIN\_FLOAT} и проверялись места её применения в алгоритме.

В момент остановки на примере \texttt{belgium\_osm} фронт (множество активных вершин) содержит 404~вершины, однако буферы фронта (массивы индексов и значений разреженного вектора) заполнены нулевыми байтами при сохранившемся значении счётчика размера фронта (числа его элементов) 404 с предыдущей итерации. Эффект воспроизводился в независимых запусках без отличий. На \texttt{roadNet-CA} картина та же: остановка на 533-й итерации при фронте из 442~вершин.

Такое сочетание соответствует незапущенному вычислительному ядру и устаревшему значению счётчика из пула повторно используемых счётчиков библиотеки (пул не обнуляет счётчики при возврате). Незапущенным оказалось вычислительное ядро \texttt{reduce\_by\_key\_small} (свёртка пар \enquote{ключ-значение} с одинаковыми ключами) с одной рабочей группой (группа рабочих элементов OpenCL)~--- одно из мест применения операции \texttt{MIN\_FLOAT} в алгоритме.

При значении переменной \texttt{prods} (число элементов для обработки) в запуске на \texttt{belgium\_osm}, равном 503, размер рабочей группы составляет 512~рабочих элементов: группа вычисляется округлением числа элементов вверх до кратного 32 (шаг wavefront~--- минимального синхронно исполняемого набора потоков данного GPU): любое значение из 481--512 даёт ту же группу 512, что и 503. У OpenCL-устройства есть общий максимум размера рабочей группы (\texttt{CL\_DEVICE\_MAX\_WORK\_GROUP\_SIZE}), а у каждого скомпилированного ядра~--- индивидуальный лимит (\texttt{CL\_KERNEL\_WORK\_GROUP\_SIZE}), запрашиваемый вызовом \texttt{clGetKernelWorkGroupInfo}, \foreignquote{english}{the maximum work-group size that can be used to execute the kernel} (максимальный размер рабочей группы, который может быть использован для запуска ядра)\footnote{URL: \url{https://manpages.debian.org/unstable/opencl-1.2-man-doc/clGetKernelWorkGroupInfo.3clc.en.html} (дата обращения:~30.09.2026).}. Запуском изолированной копии ядра \texttt{reduce\_by\_key\_small} с варьируемым размером группы и проверкой кода возврата установлено, что группы размером до 480~элементов завершаются с кодом \texttt{CL\_SUCCESS} (код успешного завершения), а группа 512 для данного ядра отклоняется драйвером с ошибкой \texttt{CL\_OUT\_OF\_RESOURCES} (недостаточно ресурсов для запуска); при этом запрос вернул для данного ядра лимит 480 при заявленном общем максимуме устройства 512. Сборка того же ядра с \texttt{fmin} получает лимит 512.

Такое поведение соответствует описанию в спецификации OpenCL~1.2~--- версии языка, под которую собраны ядра: вызов с размером вплоть до максимума устройства формально допустим, а отказ \foreignquote{english}{\texttt{CL\_OUT\_OF\_RESOURCES} \dots because of insufficient resources \dots such as registers or local memory} (недостаточно ресурсов, таких как регистры или локальная память)\footnote{URL: \url{https://manpages.debian.org/unstable/opencl-1.2-man-doc/clEnqueueNDRangeKernel.3clc.en.html} (дата обращения:~30.09.2026).}~--- штатный ответ среды исполнения. То, что индивидуальный лимит ядра может быть меньше общего максимума устройства, прямо следует из спецификации. Недочётом реализации являлось отсутствие сверки запускаемых групп с индивидуальным лимитом ядра. Код возврата постановки ядра в очередь не проверялся, поэтому программа считала запуск состоявшимся: легальный отказ драйвера выявлялся лишь внешней сверкой с эталоном.

Перед запуском вычислительных ядер с одной рабочей группой разработанное исправление сверяет размер группы с индивидуальным лимитом скомпилированного ядра. Проверяются ядро редукции \texttt{reduce\_by\_key\_small} и ядра битонной сортировки~--- параллельного алгоритма упорядочивания пар \enquote{ключ-значение} по ключу: \texttt{bitonic\_sort\_local} и \texttt{bitonic\_sort\_global}. Вспомогательная функция добавлена в \texttt{cl\_program\_builder}~--- модуль сборки OpenCL-программ. Размеры групп сверх лимита направляются по существующим веткам вычислений на малых группах (поэлементный путь через offsets+scan+scalar~--- префиксные суммы смещений групп и поэлементная редукция каждой группы, и radix-сортировка~--- поразрядная сортировка). Решение не содержит аппаратно-зависимых констант и ветвлений по производителю: лимит ядра запрашивается у драйвера.

Реализация решения осуществлялась с использованием языка C++17 и прикладного программного интерфейса OpenCL~1.2 C, выполнялась кросс-компиляция под архитектуру RISC-V~--- сборка исполняемых файлов платы на хосте (компьютере разработчика) с другой архитектурой~--- набором инструментов Syntacore\footnote{URL: \url{https://syntacore.com/tools/development-tools} (дата обращения:~01.10.2026).}, запуски проводились непосредственно на плате Milk-V Jupiter. Разработка велась в форке репозитория Spla\footnote{URL: \url{https://github.com/Redkin-Aleksey/spla} (дата обращения:~07.10.2026).}; выявленный отказ зарегистрирован как дефект\footnote{URL: \url{https://github.com/SparseLinearAlgebra/spla/issues/246} (дата обращения:~07.10.2026).}, а разработанное исправление предложено запросом на слияние\footnote{URL: \url{https://github.com/SparseLinearAlgebra/spla/pull/247} (дата обращения:~07.10.2026).}, который на момент написания отчёта находится на рассмотрении разработчиков проекта.

\section{Тестирование}

\textbf{Среда выполнения.} Измерения проводились на плате Milk-V Jupiter под ОС Bianbu 3.0.1 (Linux 6.6.63, riscv64): графический процессор \texttt{PowerVR B-Series BXE-2-32} (Imagination Technologies), драйвер версии 24.2@6603887, профиль \texttt{EMBEDDED\_PROFILE} (профиль OpenCL для встраиваемых систем), OpenCL~3.0, ядра собраны с флагом \texttt{-cl-std=CL1.2}. За основу взята ветка \texttt{main} исходного репозитория Spla на коммите \texttt{b81d1e29}. Заявленный максимум устройства \texttt{CL\_DEVICE\_MAX\_WORK\_GROUP\_SIZE = 512}, измеренный лимит отказавшего ядра \texttt{CL\_KERNEL\_WORK\_GROUP\_SIZE = 480}. Поскольку оценивается функциональная корректность, а не время выполнения, характеристики хоста на результаты вычислений не влияют: вычислительные ядра исполнялись на графическом процессоре платы, а независимая эталонная реализация~--- на её центральном процессоре; хост использовался лишь для кросс-компиляции.

\textbf{Входные данные.} В запусках использовались симметризованные версии графов из набора graph-bench\footnote{URL: \url{https://github.com/SparseLinearAlgebra/graph-bench} (дата обращения:~30.09.2026).}~--- коллекции графов и скриптов запуска для тестирования графовых библиотек~--- с единичными весами рёбер: загрузчик матриц \texttt{MtxLoader} добавляет обратные рёбра, а пример \texttt{sssp} присваивает всем рёбрам вес 1.0. Исходные матрицы в формате Matrix Market (текстовый формат для разреженных матриц)\footnote{URL: \url{https://math.nist.gov/MatrixMarket/formats.html} (дата обращения:~10.10.2026).} взяты из открытой коллекции SuiteSparse Matrix Collection\footnote{URL: \url{https://sparse.tamu.edu} (дата обращения:~07.10.2026).}; вычисляемые расстояния~--- числа с плавающей точкой одинарной точности. Характеристики использованных графов приведены в таблице~\ref{tab:graphs} для загруженных (симметризованных) версий~--- тех данных, которые непосредственно обрабатывал алгоритм; значения округлены до 0,1~млн. Граф \texttt{amazon-2008} использован как контрольный, на котором отказ не воспроизводился; \texttt{belgium\_osm} и \texttt{roadNet-CA}~--- проблемные графы, на которых отказ наблюдался.

\begin{table}[htbp]
    \centering
    \caption{Используемые наборы данных}
    \label{tab:graphs}
    \begin{tabular}{lrrl}
        \toprule
        \textbf{Граф} & \textbf{Вершины, млн} & \textbf{Рёбра, млн} & \textbf{Домен} \\
        \midrule
        \texttt{belgium\_osm} & $\approx$1,4 & $\approx$3,1 & дорожная сеть \\
        \texttt{roadNet-CA} & $\approx$2,0 & $\approx$5,5 & дорожная сеть \\
        \texttt{amazon-2008} & $\approx$0,7 & $\approx$7,0 & сеть совместных покупок \\
        \bottomrule
    \end{tabular}
\end{table}

\textbf{Методика проверки.} Целью было проверить гипотезу о том, что ограничение запускаемых групп индивидуальным лимитом ядра восстанавливает сходимость алгоритма до эталонной. Программа \texttt{sssp} из репозитория Spla сверяет CPU- и GPU-варианты с независимым эталоном (последовательная реализация без OpenCL): строка \texttt{CHECK} отмечает проверяемый вариант, строки \texttt{VERIFY} перечисляют несовпавшие вершины; отсутствие строк \texttt{VERIFY} при \texttt{CHECK acc} (проверка ускоренного варианта) означает совпадение с эталоном. Сравнивались исходная версия и версия с исправлением на тех же графах. Модель вычислений детерминирована, поэтому результаты запусков воспроизводимы без статистического разброса, отдельная статистическая обработка не проводилась.

Результаты сопоставления с эталоном представлены в таблице~\ref{tab:verify}: столбцы \enquote{До/После, расхождений}~--- число вершин, не совпавших с эталоном, \enquote{До, \% вершин}~--- их доля от всех вершин графа.

\begin{table}[htbp]
    \centering
    \caption{Сверка с эталоном до и после исправления}
    \label{tab:verify}
    \begin{tabular}{lrrr}
        \toprule
        \textbf{Граф} & \textbf{До, расхождений} & \textbf{До, \% вершин} & \textbf{После, расхождений} \\
        \midrule
        \texttt{belgium\_osm} & 1\,436\,362 & 99,7 & 0 \\
        \texttt{roadNet-CA} & 2\,155 & 0,1 & 0 \\
        \texttt{amazon-2008} & 0 & 0,0 & 0 \\
        \bottomrule
    \end{tabular}
\end{table}

До исправления первые расхождения на примере \texttt{belgium\_osm} имели вид \texttt{expected 465 actual 3.40282e+38} (ожидаемое расстояние 465~рёбер при единичных весах; значение \texttt{FLT\_MAX}: вершины недостижимы из-за преждевременного останова алгоритма). После исправления зафиксирована полная сходимость результатов на всех трёх графах: расхождений с эталоном нет.

Таким образом, на проверенных графах и конфигурациях детерминированный отказ вычислений устранён: предельный размер группы запрашивается у драйвера устройства, а корректность подтверждена совпадением с эталоном на контрольных графах.

\section{Заключение}

В ходе выполнения работы получены следующие результаты:
\begin{itemize}
    \item выполнен обзор итеративной схемы SSSP в Spla; по его итогам выбраны точки измерения для локализации;
    \item локализовано место отказа на Milk-V Jupiter и экспериментально подтверждена корневая причина: при лимите ядра 480 запускалось ядро \texttt{reduce\_by\_key\_small} с группой 512, а код \texttt{CL\_OUT\_OF\_RESOURCES} не проверялся, что затрудняло диагностику;
    \item реализовано переносимое исправление маршрутизации рабочих групп в библиотеке Spla;
    \item проведена проверка сверкой с эталоном: расхождения устранены полностью, на всех контрольных графах ноль несовпадений.
\end{itemize}

Код доступен в форке репозитория (ссылка приведена выше); дефект зарегистрирован в системе отслеживания ошибок проекта, исправление предложено запросом на слияние (ссылки приведены выше).

\end{document}
